\chapter{Resultados}
\label{cap:resultados}

Este capítulo apresenta os resultados obtidos pelas abordagens exata e heurística nos quatro portes de instância, segundo o protocolo da \autoref{tab:protocolo-experimental}. A \autoref{sec:res-comparacao} compara o Gurobi e o Algoritmo Memético em cada porte quanto ao valor da função objetivo $Z$ (\autoref{eq:objetivo}), à viabilidade das soluções heurísticas e ao esforço computacional. A \autoref{sec:res-operacionais} discute as métricas operacionais das soluções, isto é, o que cada alocação representa para o hospital em termos de especialidade, transferências, quartos mistos e ocupação. Por fim, a \autoref{sec:res-sintese} reúne as principais conclusões do capítulo.

% =====================================================================
\section{Comparação entre o Gurobi e a heurística}
\label{sec:res-comparacao}

Para cada porte, apresenta-se primeiro uma tabela com os resultados dos 10 conjuntos e, em seguida, uma figura por conjunto. Nas tabelas, as colunas do Gurobi trazem a melhor solução inteira encontrada dentro do teto de tempo (limite superior), o limite inferior do \textit{branch-and-bound} e o gap de otimalidade entre eles. Nas colunas da heurística, a melhor execução é a de menor \textit{fitness} $F$ (\autoref{eq:fitness}), e o valor reportado é o seu $Z$, sem a penalidade de capacidade, de modo que as duas abordagens são sempre comparadas pela mesma função objetivo. A média e o desvio padrão ($\sigma$) resumem as execuções independentes de cada conjunto e medem a estabilidade da heurística \cite{talbi2009metaheuristics}.

Cada figura reúne três gráficos de um conjunto: (a) a convergência da heurística, com o melhor \textit{fitness} encontrado até cada geração em todas as execuções e a melhor delas em destaque; (b) o $Z$ do Gurobi ao lado do $Z$ de cada execução da heurística e da melhor solução viável encontrada, ordenados em ordem crescente; e (c) a dispersão de $Z$ entre as execuções.

Em nenhum dos 40 conjuntos o Gurobi comprovou a otimalidade de sua solução dentro do teto de tempo. Ainda assim, quando não é trivial, o limite inferior $LB$ permite avaliar a qualidade da heurística: como nenhuma solução pode ter $Z$ menor que $LB$, a diferença relativa $(Z - LB)/Z$ indica o quanto, no máximo, a solução heurística está acima do ótimo.

% ---------------------------------------------------------------------
\subsection{Instância pequena}
\label{sec:res-pequena}

Na instância pequena (100 pacientes, 30 segundos por execução), o Gurobi obteve a melhor solução em todos os 10 conjuntos (\autoref{tab:res-z-pequeno}). A melhor execução da heurística ficou, em média, 2,0\% acima do limite superior do Gurobi (de 0,5\% a 4,9\%, conforme o conjunto), e a média das execuções, 5,6\% acima. Mesmo sem comprovar a otimalidade, o Gurobi terminou com gaps de 1,0\% a 12,3\% (5,4\% em média), o que mostra que, nesse porte, o modelo exato ainda é tratável.

Os limites inferiores também certificam a qualidade da heurística: a melhor execução está, em média, a no máximo 7,3\% do ótimo (de 2,3\% a 14,7\%). A heurística mostrou-se estável, com desvio padrão entre execuções de cerca de 2\% da média, e todas as 100 execuções terminaram com soluções viáveis.

As figuras de convergência (\autoref{fig:res-conjunto1-pequeno} a \autoref{fig:res-conjunto10-pequeno}) mostram um padrão comum aos 10 conjuntos. A maior parte da redução do \textit{fitness} ocorre na primeira geração (95\% da melhora total, em média), quando a busca local refina a população inicial, e a última melhora ocorre, em média, na 10ª das cerca de 25 gerações. Nas aproximadamente 15 gerações restantes, a melhor solução não muda, ou seja, a heurística estagna antes do fim do teto, e o tempo adicional não se converte em melhora. Nesse porte, portanto, a limitação da heurística está na capacidade de escapar dos ótimos locais, e não no tempo disponível.

\begin{table}[htbp]
    \caption{Função objetivo $Z$ por conjunto --- instância pequena. Na heurística, $Z$ sem a penalidade de capacidade; $\sigma$: desvio padrão de $Z$ entre as 10 execuções do conjunto}
    \label{tab:res-z-pequeno}
    \centering
    \begin{adjustbox}{max width=\textwidth}
    \begin{tabular}{crrrrr}
        \hline
        & \multicolumn{3}{c}{\textbf{Gurobi}} & \multicolumn{2}{c}{\textbf{Heurística (10 execuções)}} \\
        \cline{2-4} \cline{5-6}
        \textbf{Conjunto} & \textbf{Limite superior} & \textbf{Limite inferior} & \textbf{Gap} & \textbf{Melhor execução} & \textbf{Média $\pm$ $\sigma$} \\
        \hline
        01 & 11\,800,0 & 11\,050,0 & 6,4\% & 11\,975,0 & 12\,482,5 $\pm$ 304,6 \\
        02 & 11\,975,0 & 11\,850,0 & 1,0\% & 12\,125,0 & 12\,532,5 $\pm$ 232,8 \\
        03 & 9\,050,0 & 8\,125,0 & 10,2\% & 9\,150,0 & 9\,622,5 $\pm$ 219,0 \\
        04 & 5\,775,0 & 5\,600,0 & 3,0\% & 5\,900,0 & 6\,015,0 $\pm$ 102,9 \\
        05 & 8\,525,0 & 8\,025,0 & 5,9\% & 8\,775,0 & 9\,137,5 $\pm$ 267,3 \\
        06 & 16\,250,0 & 15\,350,0 & 5,5\% & 16\,375,0 & 16\,530,0 $\pm$ 107,9 \\
        07 & 5\,475,0 & 4\,800,0 & 12,3\% & 5\,625,0 & 6\,105,0 $\pm$ 246,9 \\
        08 & 11\,125,0 & 10\,600,0 & 4,7\% & 11\,675,0 & 11\,970,0 $\pm$ 220,1 \\
        09 & 9\,275,0 & 9\,100,0 & 1,9\% & 9\,325,0 & 9\,565,0 $\pm$ 209,2 \\
        10 & 14\,450,0 & 13\,950,0 & 3,5\% & 14\,750,0 & 15\,045,0 $\pm$ 143,3 \\
        \hline
    \end{tabular}
    \end{adjustbox}
\end{table}

\begin{figure}[htbp]
    \centering
    \begin{minipage}[b]{0.48\textwidth}
        \centering
        \includegraphics[width=\textwidth]{Figuras/resultados/pequeno/set_01/convergence_broken_axis.png}\\
        \small (a) Convergência da heurística
    \end{minipage}
    \hfill
    \begin{minipage}[b]{0.48\textwidth}
        \centering
        \includegraphics[width=\textwidth]{Figuras/resultados/pequeno/set_01/gurobi_vs_iterations.png}\\
        \small (b) Gurobi vs.\ execuções da heurística
    \end{minipage}

    \vspace{0.4cm}
    \begin{minipage}[b]{0.48\textwidth}
        \centering
        \includegraphics[width=\textwidth]{Figuras/resultados/pequeno/set_01/boxplot.png}\\
        \small (c) Dispersão de $Z$ entre execuções
    \end{minipage}
    \caption{Conjunto 1 --- instância pequena}
    \label{fig:res-conjunto1-pequeno}
\end{figure}
\clearpage

% Conjuntos 2 a 10: desativados até as imagens estarem no Overleaf (remover \iffalse/\fi para reativar)
\iffalse
\begin{figure}[htbp]
    \centering
    \begin{minipage}[b]{0.48\textwidth}
        \centering
        \includegraphics[width=\textwidth]{Figuras/resultados/pequeno/set_02/convergence_broken_axis.png}\\
        \small (a) Convergência da heurística
    \end{minipage}
    \hfill
    \begin{minipage}[b]{0.48\textwidth}
        \centering
        \includegraphics[width=\textwidth]{Figuras/resultados/pequeno/set_02/gurobi_vs_iterations.png}\\
        \small (b) Gurobi vs.\ execuções da heurística
    \end{minipage}

    \vspace{0.4cm}
    \begin{minipage}[b]{0.48\textwidth}
        \centering
        \includegraphics[width=\textwidth]{Figuras/resultados/pequeno/set_02/boxplot.png}\\
        \small (c) Dispersão de $Z$ entre execuções
    \end{minipage}
    \caption{Conjunto 2 --- instância pequena}
    \label{fig:res-conjunto2-pequeno}
\end{figure}
\clearpage

\begin{figure}[htbp]
    \centering
    \begin{minipage}[b]{0.48\textwidth}
        \centering
        \includegraphics[width=\textwidth]{Figuras/resultados/pequeno/set_03/convergence_broken_axis.png}\\
        \small (a) Convergência da heurística
    \end{minipage}
    \hfill
    \begin{minipage}[b]{0.48\textwidth}
        \centering
        \includegraphics[width=\textwidth]{Figuras/resultados/pequeno/set_03/gurobi_vs_iterations.png}\\
        \small (b) Gurobi vs.\ execuções da heurística
    \end{minipage}

    \vspace{0.4cm}
    \begin{minipage}[b]{0.48\textwidth}
        \centering
        \includegraphics[width=\textwidth]{Figuras/resultados/pequeno/set_03/boxplot.png}\\
        \small (c) Dispersão de $Z$ entre execuções
    \end{minipage}
    \caption{Conjunto 3 --- instância pequena}
    \label{fig:res-conjunto3-pequeno}
\end{figure}
\clearpage

\begin{figure}[htbp]
    \centering
    \begin{minipage}[b]{0.48\textwidth}
        \centering
        \includegraphics[width=\textwidth]{Figuras/resultados/pequeno/set_04/convergence_broken_axis.png}\\
        \small (a) Convergência da heurística
    \end{minipage}
    \hfill
    \begin{minipage}[b]{0.48\textwidth}
        \centering
        \includegraphics[width=\textwidth]{Figuras/resultados/pequeno/set_04/gurobi_vs_iterations.png}\\
        \small (b) Gurobi vs.\ execuções da heurística
    \end{minipage}

    \vspace{0.4cm}
    \begin{minipage}[b]{0.48\textwidth}
        \centering
        \includegraphics[width=\textwidth]{Figuras/resultados/pequeno/set_04/boxplot.png}\\
        \small (c) Dispersão de $Z$ entre execuções
    \end{minipage}
    \caption{Conjunto 4 --- instância pequena}
    \label{fig:res-conjunto4-pequeno}
\end{figure}
\clearpage

\begin{figure}[htbp]
    \centering
    \begin{minipage}[b]{0.48\textwidth}
        \centering
        \includegraphics[width=\textwidth]{Figuras/resultados/pequeno/set_05/convergence_broken_axis.png}\\
        \small (a) Convergência da heurística
    \end{minipage}
    \hfill
    \begin{minipage}[b]{0.48\textwidth}
        \centering
        \includegraphics[width=\textwidth]{Figuras/resultados/pequeno/set_05/gurobi_vs_iterations.png}\\
        \small (b) Gurobi vs.\ execuções da heurística
    \end{minipage}

    \vspace{0.4cm}
    \begin{minipage}[b]{0.48\textwidth}
        \centering
        \includegraphics[width=\textwidth]{Figuras/resultados/pequeno/set_05/boxplot.png}\\
        \small (c) Dispersão de $Z$ entre execuções
    \end{minipage}
    \caption{Conjunto 5 --- instância pequena}
    \label{fig:res-conjunto5-pequeno}
\end{figure}
\clearpage

\begin{figure}[htbp]
    \centering
    \begin{minipage}[b]{0.48\textwidth}
        \centering
        \includegraphics[width=\textwidth]{Figuras/resultados/pequeno/set_06/convergence_broken_axis.png}\\
        \small (a) Convergência da heurística
    \end{minipage}
    \hfill
    \begin{minipage}[b]{0.48\textwidth}
        \centering
        \includegraphics[width=\textwidth]{Figuras/resultados/pequeno/set_06/gurobi_vs_iterations.png}\\
        \small (b) Gurobi vs.\ execuções da heurística
    \end{minipage}

    \vspace{0.4cm}
    \begin{minipage}[b]{0.48\textwidth}
        \centering
        \includegraphics[width=\textwidth]{Figuras/resultados/pequeno/set_06/boxplot.png}\\
        \small (c) Dispersão de $Z$ entre execuções
    \end{minipage}
    \caption{Conjunto 6 --- instância pequena}
    \label{fig:res-conjunto6-pequeno}
\end{figure}
\clearpage

\begin{figure}[htbp]
    \centering
    \begin{minipage}[b]{0.48\textwidth}
        \centering
        \includegraphics[width=\textwidth]{Figuras/resultados/pequeno/set_07/convergence_broken_axis.png}\\
        \small (a) Convergência da heurística
    \end{minipage}
    \hfill
    \begin{minipage}[b]{0.48\textwidth}
        \centering
        \includegraphics[width=\textwidth]{Figuras/resultados/pequeno/set_07/gurobi_vs_iterations.png}\\
        \small (b) Gurobi vs.\ execuções da heurística
    \end{minipage}

    \vspace{0.4cm}
    \begin{minipage}[b]{0.48\textwidth}
        \centering
        \includegraphics[width=\textwidth]{Figuras/resultados/pequeno/set_07/boxplot.png}\\
        \small (c) Dispersão de $Z$ entre execuções
    \end{minipage}
    \caption{Conjunto 7 --- instância pequena}
    \label{fig:res-conjunto7-pequeno}
\end{figure}
\clearpage

\begin{figure}[htbp]
    \centering
    \begin{minipage}[b]{0.48\textwidth}
        \centering
        \includegraphics[width=\textwidth]{Figuras/resultados/pequeno/set_08/convergence_broken_axis.png}\\
        \small (a) Convergência da heurística
    \end{minipage}
    \hfill
    \begin{minipage}[b]{0.48\textwidth}
        \centering
        \includegraphics[width=\textwidth]{Figuras/resultados/pequeno/set_08/gurobi_vs_iterations.png}\\
        \small (b) Gurobi vs.\ execuções da heurística
    \end{minipage}

    \vspace{0.4cm}
    \begin{minipage}[b]{0.48\textwidth}
        \centering
        \includegraphics[width=\textwidth]{Figuras/resultados/pequeno/set_08/boxplot.png}\\
        \small (c) Dispersão de $Z$ entre execuções
    \end{minipage}
    \caption{Conjunto 8 --- instância pequena}
    \label{fig:res-conjunto8-pequeno}
\end{figure}
\clearpage

\begin{figure}[htbp]
    \centering
    \begin{minipage}[b]{0.48\textwidth}
        \centering
        \includegraphics[width=\textwidth]{Figuras/resultados/pequeno/set_09/convergence_broken_axis.png}\\
        \small (a) Convergência da heurística
    \end{minipage}
    \hfill
    \begin{minipage}[b]{0.48\textwidth}
        \centering
        \includegraphics[width=\textwidth]{Figuras/resultados/pequeno/set_09/gurobi_vs_iterations.png}\\
        \small (b) Gurobi vs.\ execuções da heurística
    \end{minipage}

    \vspace{0.4cm}
    \begin{minipage}[b]{0.48\textwidth}
        \centering
        \includegraphics[width=\textwidth]{Figuras/resultados/pequeno/set_09/boxplot.png}\\
        \small (c) Dispersão de $Z$ entre execuções
    \end{minipage}
    \caption{Conjunto 9 --- instância pequena}
    \label{fig:res-conjunto9-pequeno}
\end{figure}
\clearpage

\begin{figure}[htbp]
    \centering
    \begin{minipage}[b]{0.48\textwidth}
        \centering
        \includegraphics[width=\textwidth]{Figuras/resultados/pequeno/set_10/convergence_broken_axis.png}\\
        \small (a) Convergência da heurística
    \end{minipage}
    \hfill
    \begin{minipage}[b]{0.48\textwidth}
        \centering
        \includegraphics[width=\textwidth]{Figuras/resultados/pequeno/set_10/gurobi_vs_iterations.png}\\
        \small (b) Gurobi vs.\ execuções da heurística
    \end{minipage}

    \vspace{0.4cm}
    \begin{minipage}[b]{0.48\textwidth}
        \centering
        \includegraphics[width=\textwidth]{Figuras/resultados/pequeno/set_10/boxplot.png}\\
        \small (c) Dispersão de $Z$ entre execuções
    \end{minipage}
    \caption{Conjunto 10 --- instância pequena}
    \label{fig:res-conjunto10-pequeno}
\end{figure}
\clearpage
\fi

% ---------------------------------------------------------------------
\subsection{Instância média}
\label{sec:res-media}

Na instância média (400 pacientes, 90 segundos por execução), ocorre a inversão entre as abordagens: a heurística obteve o menor $Z$ em todos os 10 conjuntos (\autoref{tab:res-z-medio}), em média 65,5\% abaixo do limite superior do Gurobi (de 54,3\% a 77,3\%). A vantagem não depende de uma execução favorável, pois a média das execuções também ficou 64,4\% abaixo do Gurobi.

Com quatro vezes mais pacientes que a instância pequena, o Gurobi terminou o teto com gaps de 60,5\% a 81,3\%, ainda distante tanto de uma boa solução quanto de um limite inferior justo. Esses limites, porém, já não são triviais e confirmam a qualidade da heurística: a melhor execução está, em média, a no máximo 17,2\% do ótimo (de 12,5\% a 23,9\%). Como o próprio limite inferior tende a estar abaixo do ótimo, a distância real é provavelmente menor.

Todas as 100 execuções encontraram ao menos uma solução viável e, em 99 delas, a melhor solução da execução também era viável. O desvio padrão entre execuções permaneceu em torno de 2\% da média. Ao contrário da instância pequena, as curvas de convergência (\autoref{fig:res-conjunto1-medio} a \autoref{fig:res-conjunto10-medio}) ainda descem ao final do teto: a última melhora ocorre, em média, na geração 14,7 de 15,5. Nesse porte, a heurística é limitada pelo tempo disponível, e não pela estagnação.

\begin{table}[htbp]
    \caption{Função objetivo $Z$ por conjunto --- instância média. Na heurística, $Z$ sem a penalidade de capacidade; $\sigma$: desvio padrão de $Z$ entre as 10 execuções do conjunto}
    \label{tab:res-z-medio}
    \centering
    \begin{adjustbox}{max width=\textwidth}
    \begin{tabular}{crrrrr}
        \hline
        & \multicolumn{3}{c}{\textbf{Gurobi}} & \multicolumn{2}{c}{\textbf{Heurística (10 execuções)}} \\
        \cline{2-4} \cline{5-6}
        \textbf{Conjunto} & \textbf{Limite superior} & \textbf{Limite inferior} & \textbf{Gap} & \textbf{Melhor execução} & \textbf{Média $\pm$ $\sigma$} \\
        \hline
        01 & 79\,300,0 & 22\,350,0 & 71,8\% & 29\,350,0 & 30\,077,5 $\pm$ 585,8 \\
        02 & 73\,625,0 & 18\,000,0 & 75,6\% & 22\,750,0 & 23\,655,0 $\pm$ 545,2 \\
        03 & 91\,800,0 & 27\,400,0 & 70,2\% & 31\,675,0 & 32\,997,5 $\pm$ 711,4 \\
        04 & 84\,525,0 & 23\,800,0 & 71,8\% & 28\,475,0 & 29\,135,0 $\pm$ 386,3 \\
        05 & 89\,950,0 & 20\,450,0 & 77,3\% & 25\,150,0 & 26\,827,5 $\pm$ 875,6 \\
        06 & 69\,075,0 & 24\,550,0 & 64,5\% & 29\,575,0 & 30\,532,5 $\pm$ 578,8 \\
        07 & 83\,950,0 & 22\,750,0 & 72,9\% & 27\,850,0 & 28\,812,5 $\pm$ 765,0 \\
        08 & 88\,275,0 & 16\,500,0 & 81,3\% & 20\,000,0 & 21\,027,5 $\pm$ 651,3 \\
        09 & 92\,650,0 & 29\,400,0 & 68,3\% & 33\,600,0 & 34\,217,5 $\pm$ 598,2 \\
        10 & 86\,775,0 & 34\,300,0 & 60,5\% & 39\,650,0 & 40\,195,0 $\pm$ 384,9 \\
        \hline
    \end{tabular}
    \end{adjustbox}
\end{table}

\begin{figure}[htbp]
    \centering
    \begin{minipage}[b]{0.48\textwidth}
        \centering
        \includegraphics[width=\textwidth]{Figuras/resultados/medio/set_01/convergence_broken_axis.png}\\
        \small (a) Convergência da heurística
    \end{minipage}
    \hfill
    \begin{minipage}[b]{0.48\textwidth}
        \centering
        \includegraphics[width=\textwidth]{Figuras/resultados/medio/set_01/gurobi_vs_iterations.png}\\
        \small (b) Gurobi vs.\ execuções da heurística
    \end{minipage}

    \vspace{0.4cm}
    \begin{minipage}[b]{0.48\textwidth}
        \centering
        \includegraphics[width=\textwidth]{Figuras/resultados/medio/set_01/boxplot.png}\\
        \small (c) Dispersão de $Z$ entre execuções
    \end{minipage}
    \caption{Conjunto 1 --- instância média}
    \label{fig:res-conjunto1-medio}
\end{figure}
\clearpage

% Conjuntos 2 a 10: desativados até as imagens estarem no Overleaf (remover \iffalse/\fi para reativar)
\iffalse
\begin{figure}[htbp]
    \centering
    \begin{minipage}[b]{0.48\textwidth}
        \centering
        \includegraphics[width=\textwidth]{Figuras/resultados/medio/set_02/convergence_broken_axis.png}\\
        \small (a) Convergência da heurística
    \end{minipage}
    \hfill
    \begin{minipage}[b]{0.48\textwidth}
        \centering
        \includegraphics[width=\textwidth]{Figuras/resultados/medio/set_02/gurobi_vs_iterations.png}\\
        \small (b) Gurobi vs.\ execuções da heurística
    \end{minipage}

    \vspace{0.4cm}
    \begin{minipage}[b]{0.48\textwidth}
        \centering
        \includegraphics[width=\textwidth]{Figuras/resultados/medio/set_02/boxplot.png}\\
        \small (c) Dispersão de $Z$ entre execuções
    \end{minipage}
    \caption{Conjunto 2 --- instância média}
    \label{fig:res-conjunto2-medio}
\end{figure}
\clearpage

\begin{figure}[htbp]
    \centering
    \begin{minipage}[b]{0.48\textwidth}
        \centering
        \includegraphics[width=\textwidth]{Figuras/resultados/medio/set_03/convergence_broken_axis.png}\\
        \small (a) Convergência da heurística
    \end{minipage}
    \hfill
    \begin{minipage}[b]{0.48\textwidth}
        \centering
        \includegraphics[width=\textwidth]{Figuras/resultados/medio/set_03/gurobi_vs_iterations.png}\\
        \small (b) Gurobi vs.\ execuções da heurística
    \end{minipage}

    \vspace{0.4cm}
    \begin{minipage}[b]{0.48\textwidth}
        \centering
        \includegraphics[width=\textwidth]{Figuras/resultados/medio/set_03/boxplot.png}\\
        \small (c) Dispersão de $Z$ entre execuções
    \end{minipage}
    \caption{Conjunto 3 --- instância média}
    \label{fig:res-conjunto3-medio}
\end{figure}
\clearpage

\begin{figure}[htbp]
    \centering
    \begin{minipage}[b]{0.48\textwidth}
        \centering
        \includegraphics[width=\textwidth]{Figuras/resultados/medio/set_04/convergence_broken_axis.png}\\
        \small (a) Convergência da heurística
    \end{minipage}
    \hfill
    \begin{minipage}[b]{0.48\textwidth}
        \centering
        \includegraphics[width=\textwidth]{Figuras/resultados/medio/set_04/gurobi_vs_iterations.png}\\
        \small (b) Gurobi vs.\ execuções da heurística
    \end{minipage}

    \vspace{0.4cm}
    \begin{minipage}[b]{0.48\textwidth}
        \centering
        \includegraphics[width=\textwidth]{Figuras/resultados/medio/set_04/boxplot.png}\\
        \small (c) Dispersão de $Z$ entre execuções
    \end{minipage}
    \caption{Conjunto 4 --- instância média}
    \label{fig:res-conjunto4-medio}
\end{figure}
\clearpage

\begin{figure}[htbp]
    \centering
    \begin{minipage}[b]{0.48\textwidth}
        \centering
        \includegraphics[width=\textwidth]{Figuras/resultados/medio/set_05/convergence_broken_axis.png}\\
        \small (a) Convergência da heurística
    \end{minipage}
    \hfill
    \begin{minipage}[b]{0.48\textwidth}
        \centering
        \includegraphics[width=\textwidth]{Figuras/resultados/medio/set_05/gurobi_vs_iterations.png}\\
        \small (b) Gurobi vs.\ execuções da heurística
    \end{minipage}

    \vspace{0.4cm}
    \begin{minipage}[b]{0.48\textwidth}
        \centering
        \includegraphics[width=\textwidth]{Figuras/resultados/medio/set_05/boxplot.png}\\
        \small (c) Dispersão de $Z$ entre execuções
    \end{minipage}
    \caption{Conjunto 5 --- instância média}
    \label{fig:res-conjunto5-medio}
\end{figure}
\clearpage

\begin{figure}[htbp]
    \centering
    \begin{minipage}[b]{0.48\textwidth}
        \centering
        \includegraphics[width=\textwidth]{Figuras/resultados/medio/set_06/convergence_broken_axis.png}\\
        \small (a) Convergência da heurística
    \end{minipage}
    \hfill
    \begin{minipage}[b]{0.48\textwidth}
        \centering
        \includegraphics[width=\textwidth]{Figuras/resultados/medio/set_06/gurobi_vs_iterations.png}\\
        \small (b) Gurobi vs.\ execuções da heurística
    \end{minipage}

    \vspace{0.4cm}
    \begin{minipage}[b]{0.48\textwidth}
        \centering
        \includegraphics[width=\textwidth]{Figuras/resultados/medio/set_06/boxplot.png}\\
        \small (c) Dispersão de $Z$ entre execuções
    \end{minipage}
    \caption{Conjunto 6 --- instância média}
    \label{fig:res-conjunto6-medio}
\end{figure}
\clearpage

\begin{figure}[htbp]
    \centering
    \begin{minipage}[b]{0.48\textwidth}
        \centering
        \includegraphics[width=\textwidth]{Figuras/resultados/medio/set_07/convergence_broken_axis.png}\\
        \small (a) Convergência da heurística
    \end{minipage}
    \hfill
    \begin{minipage}[b]{0.48\textwidth}
        \centering
        \includegraphics[width=\textwidth]{Figuras/resultados/medio/set_07/gurobi_vs_iterations.png}\\
        \small (b) Gurobi vs.\ execuções da heurística
    \end{minipage}

    \vspace{0.4cm}
    \begin{minipage}[b]{0.48\textwidth}
        \centering
        \includegraphics[width=\textwidth]{Figuras/resultados/medio/set_07/boxplot.png}\\
        \small (c) Dispersão de $Z$ entre execuções
    \end{minipage}
    \caption{Conjunto 7 --- instância média}
    \label{fig:res-conjunto7-medio}
\end{figure}
\clearpage

\begin{figure}[htbp]
    \centering
    \begin{minipage}[b]{0.48\textwidth}
        \centering
        \includegraphics[width=\textwidth]{Figuras/resultados/medio/set_08/convergence_broken_axis.png}\\
        \small (a) Convergência da heurística
    \end{minipage}
    \hfill
    \begin{minipage}[b]{0.48\textwidth}
        \centering
        \includegraphics[width=\textwidth]{Figuras/resultados/medio/set_08/gurobi_vs_iterations.png}\\
        \small (b) Gurobi vs.\ execuções da heurística
    \end{minipage}

    \vspace{0.4cm}
    \begin{minipage}[b]{0.48\textwidth}
        \centering
        \includegraphics[width=\textwidth]{Figuras/resultados/medio/set_08/boxplot.png}\\
        \small (c) Dispersão de $Z$ entre execuções
    \end{minipage}
    \caption{Conjunto 8 --- instância média}
    \label{fig:res-conjunto8-medio}
\end{figure}
\clearpage

\begin{figure}[htbp]
    \centering
    \begin{minipage}[b]{0.48\textwidth}
        \centering
        \includegraphics[width=\textwidth]{Figuras/resultados/medio/set_09/convergence_broken_axis.png}\\
        \small (a) Convergência da heurística
    \end{minipage}
    \hfill
    \begin{minipage}[b]{0.48\textwidth}
        \centering
        \includegraphics[width=\textwidth]{Figuras/resultados/medio/set_09/gurobi_vs_iterations.png}\\
        \small (b) Gurobi vs.\ execuções da heurística
    \end{minipage}

    \vspace{0.4cm}
    \begin{minipage}[b]{0.48\textwidth}
        \centering
        \includegraphics[width=\textwidth]{Figuras/resultados/medio/set_09/boxplot.png}\\
        \small (c) Dispersão de $Z$ entre execuções
    \end{minipage}
    \caption{Conjunto 9 --- instância média}
    \label{fig:res-conjunto9-medio}
\end{figure}
\clearpage

\begin{figure}[htbp]
    \centering
    \begin{minipage}[b]{0.48\textwidth}
        \centering
        \includegraphics[width=\textwidth]{Figuras/resultados/medio/set_10/convergence_broken_axis.png}\\
        \small (a) Convergência da heurística
    \end{minipage}
    \hfill
    \begin{minipage}[b]{0.48\textwidth}
        \centering
        \includegraphics[width=\textwidth]{Figuras/resultados/medio/set_10/gurobi_vs_iterations.png}\\
        \small (b) Gurobi vs.\ execuções da heurística
    \end{minipage}

    \vspace{0.4cm}
    \begin{minipage}[b]{0.48\textwidth}
        \centering
        \includegraphics[width=\textwidth]{Figuras/resultados/medio/set_10/boxplot.png}\\
        \small (c) Dispersão de $Z$ entre execuções
    \end{minipage}
    \caption{Conjunto 10 --- instância média}
    \label{fig:res-conjunto10-medio}
\end{figure}
\clearpage
\fi

% ---------------------------------------------------------------------
\subsection{Instância grande}
\label{sec:res-grande}

Na instância grande (1600 pacientes, 360 segundos por execução), a diferença entre as abordagens aumenta: a melhor execução da heurística obteve $Z$, em média, 93,1\% menor que o do Gurobi (de 91,4\% a 94,8\%), com vantagem em todos os conjuntos (\autoref{tab:res-z-grande}). O Gurobi não avançou além da relaxação inicial em 8 dos 10 conjuntos, terminando com limite inferior igual a zero. Nos outros dois (conjuntos 3 e 6), os limites inferiores indicam que a heurística está a no máximo 32,4\% e 43,9\% do ótimo, respectivamente. Esse comportamento é coerente com a literatura: Demeester et al. \cite{demeester2010hybrid} não obtiveram sequer uma solução viável em uma hora de execução de um solver comercial em instâncias de tamanho real, e Ceschia e Schaerf \cite{ceschia2012modeling} precisaram fracionar as instâncias maiores para tornar o modelo tratável.

Nesse porte aparece a principal ressalva da heurística: a melhor solução de 7 dos 10 conjuntos viola a capacidade, ainda que de forma residual (\autoref{sec:res-viabilidade}). A violação também explica por que, nos conjuntos 1, 4, 5 e 9, o $Z$ da melhor execução é maior que a média das execuções. A melhor execução é escolhida pelo \textit{fitness} $F$, que inclui a penalidade $W_{cap} \cdot V_{cap}$, e outra execução pode ter alcançado um $Z$ menor violando mais a capacidade.

As curvas de convergência (\autoref{fig:res-conjunto1-grande} a \autoref{fig:res-conjunto10-grande}) repetem o comportamento da instância média, com melhora até as últimas gerações, e o desvio padrão entre execuções permaneceu em 2,5\% da média.

\begin{table}[htbp]
    \caption{Função objetivo $Z$ por conjunto --- instância grande. Na heurística, $Z$ sem a penalidade de capacidade; $\sigma$: desvio padrão de $Z$ entre as 10 execuções do conjunto}
    \label{tab:res-z-grande}
    \centering
    \begin{adjustbox}{max width=\textwidth}
    \begin{tabular}{crrrrr}
        \hline
        & \multicolumn{3}{c}{\textbf{Gurobi}} & \multicolumn{2}{c}{\textbf{Heurística (10 execuções)}} \\
        \cline{2-4} \cline{5-6}
        \textbf{Conjunto} & \textbf{Limite superior} & \textbf{Limite inferior} & \textbf{Gap} & \textbf{Melhor execução} & \textbf{Média $\pm$ $\sigma$} \\
        \hline
        01 & 811\,250,0 & 0,0 & 100,0\% & 56\,925,0 & 56\,180,0 $\pm$ 1\,298,9 \\
        02 & 819\,575,0 & 0,0 & 100,0\% & 68\,025,0 & 69\,017,5 $\pm$ 1\,099,0 \\
        03 & 820\,500,0 & 47\,700,0 & 94,2\% & 70\,575,0 & 72\,332,5 $\pm$ 1\,693,4 \\
        04 & 830\,875,0 & 0,0 & 100,0\% & 47\,075,0 & 46\,955,0 $\pm$ 867,4 \\
        05 & 816\,925,0 & 0,0 & 100,0\% & 60\,275,0 & 59\,227,5 $\pm$ 1\,272,4 \\
        06 & 815\,200,0 & 29\,650,0 & 96,4\% & 52\,850,0 & 53\,197,5 $\pm$ 1\,378,4 \\
        07 & 821\,300,0 & 0,0 & 100,0\% & 45\,100,0 & 45\,462,5 $\pm$ 1\,931,3 \\
        08 & 835\,450,0 & 0,0 & 100,0\% & 66\,300,0 & 67\,305,0 $\pm$ 1\,473,6 \\
        09 & 809\,600,0 & 0,0 & 100,0\% & 42\,175,0 & 41\,815,0 $\pm$ 890,0 \\
        10 & 822\,025,0 & 0,0 & 100,0\% & 53\,300,0 & 56\,097,5 $\pm$ 2\,001,9 \\
        \hline
    \end{tabular}
    \end{adjustbox}
\end{table}

\begin{figure}[htbp]
    \centering
    \begin{minipage}[b]{0.48\textwidth}
        \centering
        \includegraphics[width=\textwidth]{Figuras/resultados/grande/set_01/convergence_broken_axis.png}\\
        \small (a) Convergência da heurística
    \end{minipage}
    \hfill
    \begin{minipage}[b]{0.48\textwidth}
        \centering
        \includegraphics[width=\textwidth]{Figuras/resultados/grande/set_01/gurobi_vs_iterations.png}\\
        \small (b) Gurobi vs.\ execuções da heurística
    \end{minipage}

    \vspace{0.4cm}
    \begin{minipage}[b]{0.48\textwidth}
        \centering
        \includegraphics[width=\textwidth]{Figuras/resultados/grande/set_01/boxplot.png}\\
        \small (c) Dispersão de $Z$ entre execuções
    \end{minipage}
    \caption{Conjunto 1 --- instância grande}
    \label{fig:res-conjunto1-grande}
\end{figure}
\clearpage

% Conjuntos 2 a 10: desativados até as imagens estarem no Overleaf (remover \iffalse/\fi para reativar)
\iffalse
\begin{figure}[htbp]
    \centering
    \begin{minipage}[b]{0.48\textwidth}
        \centering
        \includegraphics[width=\textwidth]{Figuras/resultados/grande/set_02/convergence_broken_axis.png}\\
        \small (a) Convergência da heurística
    \end{minipage}
    \hfill
    \begin{minipage}[b]{0.48\textwidth}
        \centering
        \includegraphics[width=\textwidth]{Figuras/resultados/grande/set_02/gurobi_vs_iterations.png}\\
        \small (b) Gurobi vs.\ execuções da heurística
    \end{minipage}

    \vspace{0.4cm}
    \begin{minipage}[b]{0.48\textwidth}
        \centering
        \includegraphics[width=\textwidth]{Figuras/resultados/grande/set_02/boxplot.png}\\
        \small (c) Dispersão de $Z$ entre execuções
    \end{minipage}
    \caption{Conjunto 2 --- instância grande}
    \label{fig:res-conjunto2-grande}
\end{figure}
\clearpage

\begin{figure}[htbp]
    \centering
    \begin{minipage}[b]{0.48\textwidth}
        \centering
        \includegraphics[width=\textwidth]{Figuras/resultados/grande/set_03/convergence_broken_axis.png}\\
        \small (a) Convergência da heurística
    \end{minipage}
    \hfill
    \begin{minipage}[b]{0.48\textwidth}
        \centering
        \includegraphics[width=\textwidth]{Figuras/resultados/grande/set_03/gurobi_vs_iterations.png}\\
        \small (b) Gurobi vs.\ execuções da heurística
    \end{minipage}

    \vspace{0.4cm}
    \begin{minipage}[b]{0.48\textwidth}
        \centering
        \includegraphics[width=\textwidth]{Figuras/resultados/grande/set_03/boxplot.png}\\
        \small (c) Dispersão de $Z$ entre execuções
    \end{minipage}
    \caption{Conjunto 3 --- instância grande}
    \label{fig:res-conjunto3-grande}
\end{figure}
\clearpage

\begin{figure}[htbp]
    \centering
    \begin{minipage}[b]{0.48\textwidth}
        \centering
        \includegraphics[width=\textwidth]{Figuras/resultados/grande/set_04/convergence_broken_axis.png}\\
        \small (a) Convergência da heurística
    \end{minipage}
    \hfill
    \begin{minipage}[b]{0.48\textwidth}
        \centering
        \includegraphics[width=\textwidth]{Figuras/resultados/grande/set_04/gurobi_vs_iterations.png}\\
        \small (b) Gurobi vs.\ execuções da heurística
    \end{minipage}

    \vspace{0.4cm}
    \begin{minipage}[b]{0.48\textwidth}
        \centering
        \includegraphics[width=\textwidth]{Figuras/resultados/grande/set_04/boxplot.png}\\
        \small (c) Dispersão de $Z$ entre execuções
    \end{minipage}
    \caption{Conjunto 4 --- instância grande}
    \label{fig:res-conjunto4-grande}
\end{figure}
\clearpage

\begin{figure}[htbp]
    \centering
    \begin{minipage}[b]{0.48\textwidth}
        \centering
        \includegraphics[width=\textwidth]{Figuras/resultados/grande/set_05/convergence_broken_axis.png}\\
        \small (a) Convergência da heurística
    \end{minipage}
    \hfill
    \begin{minipage}[b]{0.48\textwidth}
        \centering
        \includegraphics[width=\textwidth]{Figuras/resultados/grande/set_05/gurobi_vs_iterations.png}\\
        \small (b) Gurobi vs.\ execuções da heurística
    \end{minipage}

    \vspace{0.4cm}
    \begin{minipage}[b]{0.48\textwidth}
        \centering
        \includegraphics[width=\textwidth]{Figuras/resultados/grande/set_05/boxplot.png}\\
        \small (c) Dispersão de $Z$ entre execuções
    \end{minipage}
    \caption{Conjunto 5 --- instância grande}
    \label{fig:res-conjunto5-grande}
\end{figure}
\clearpage

\begin{figure}[htbp]
    \centering
    \begin{minipage}[b]{0.48\textwidth}
        \centering
        \includegraphics[width=\textwidth]{Figuras/resultados/grande/set_06/convergence_broken_axis.png}\\
        \small (a) Convergência da heurística
    \end{minipage}
    \hfill
    \begin{minipage}[b]{0.48\textwidth}
        \centering
        \includegraphics[width=\textwidth]{Figuras/resultados/grande/set_06/gurobi_vs_iterations.png}\\
        \small (b) Gurobi vs.\ execuções da heurística
    \end{minipage}

    \vspace{0.4cm}
    \begin{minipage}[b]{0.48\textwidth}
        \centering
        \includegraphics[width=\textwidth]{Figuras/resultados/grande/set_06/boxplot.png}\\
        \small (c) Dispersão de $Z$ entre execuções
    \end{minipage}
    \caption{Conjunto 6 --- instância grande}
    \label{fig:res-conjunto6-grande}
\end{figure}
\clearpage

\begin{figure}[htbp]
    \centering
    \begin{minipage}[b]{0.48\textwidth}
        \centering
        \includegraphics[width=\textwidth]{Figuras/resultados/grande/set_07/convergence_broken_axis.png}\\
        \small (a) Convergência da heurística
    \end{minipage}
    \hfill
    \begin{minipage}[b]{0.48\textwidth}
        \centering
        \includegraphics[width=\textwidth]{Figuras/resultados/grande/set_07/gurobi_vs_iterations.png}\\
        \small (b) Gurobi vs.\ execuções da heurística
    \end{minipage}

    \vspace{0.4cm}
    \begin{minipage}[b]{0.48\textwidth}
        \centering
        \includegraphics[width=\textwidth]{Figuras/resultados/grande/set_07/boxplot.png}\\
        \small (c) Dispersão de $Z$ entre execuções
    \end{minipage}
    \caption{Conjunto 7 --- instância grande}
    \label{fig:res-conjunto7-grande}
\end{figure}
\clearpage

\begin{figure}[htbp]
    \centering
    \begin{minipage}[b]{0.48\textwidth}
        \centering
        \includegraphics[width=\textwidth]{Figuras/resultados/grande/set_08/convergence_broken_axis.png}\\
        \small (a) Convergência da heurística
    \end{minipage}
    \hfill
    \begin{minipage}[b]{0.48\textwidth}
        \centering
        \includegraphics[width=\textwidth]{Figuras/resultados/grande/set_08/gurobi_vs_iterations.png}\\
        \small (b) Gurobi vs.\ execuções da heurística
    \end{minipage}

    \vspace{0.4cm}
    \begin{minipage}[b]{0.48\textwidth}
        \centering
        \includegraphics[width=\textwidth]{Figuras/resultados/grande/set_08/boxplot.png}\\
        \small (c) Dispersão de $Z$ entre execuções
    \end{minipage}
    \caption{Conjunto 8 --- instância grande}
    \label{fig:res-conjunto8-grande}
\end{figure}
\clearpage

\begin{figure}[htbp]
    \centering
    \begin{minipage}[b]{0.48\textwidth}
        \centering
        \includegraphics[width=\textwidth]{Figuras/resultados/grande/set_09/convergence_broken_axis.png}\\
        \small (a) Convergência da heurística
    \end{minipage}
    \hfill
    \begin{minipage}[b]{0.48\textwidth}
        \centering
        \includegraphics[width=\textwidth]{Figuras/resultados/grande/set_09/gurobi_vs_iterations.png}\\
        \small (b) Gurobi vs.\ execuções da heurística
    \end{minipage}

    \vspace{0.4cm}
    \begin{minipage}[b]{0.48\textwidth}
        \centering
        \includegraphics[width=\textwidth]{Figuras/resultados/grande/set_09/boxplot.png}\\
        \small (c) Dispersão de $Z$ entre execuções
    \end{minipage}
    \caption{Conjunto 9 --- instância grande}
    \label{fig:res-conjunto9-grande}
\end{figure}
\clearpage

\begin{figure}[htbp]
    \centering
    \begin{minipage}[b]{0.48\textwidth}
        \centering
        \includegraphics[width=\textwidth]{Figuras/resultados/grande/set_10/convergence_broken_axis.png}\\
        \small (a) Convergência da heurística
    \end{minipage}
    \hfill
    \begin{minipage}[b]{0.48\textwidth}
        \centering
        \includegraphics[width=\textwidth]{Figuras/resultados/grande/set_10/gurobi_vs_iterations.png}\\
        \small (b) Gurobi vs.\ execuções da heurística
    \end{minipage}

    \vspace{0.4cm}
    \begin{minipage}[b]{0.48\textwidth}
        \centering
        \includegraphics[width=\textwidth]{Figuras/resultados/grande/set_10/boxplot.png}\\
        \small (c) Dispersão de $Z$ entre execuções
    \end{minipage}
    \caption{Conjunto 10 --- instância grande}
    \label{fig:res-conjunto10-grande}
\end{figure}
\clearpage
\fi

% ---------------------------------------------------------------------
\subsection{Instância muito grande}
\label{sec:res-muito-grande}

Na instância muito grande (6400 pacientes, 900 segundos por execução), a heurística novamente obteve o menor $Z$ em todos os conjuntos, em média 93,6\% abaixo do Gurobi (de 92,2\% a 94,4\%), como mostra a \autoref{tab:res-z-muito-grande}. O Gurobi terminou todos os conjuntos com limite inferior nulo, sem concluir a relaxação inicial, de modo que a comparação se restringe às soluções inteiras encontradas, como antecipado na \autoref{sec:complexidade}.

Nenhuma das 50 execuções encontrou solução viável, embora a violação permaneça abaixo de 0,1\% da carga de internação (\autoref{sec:res-viabilidade}). Como na instância grande, a escolha da melhor execução pelo \textit{fitness} faz com que, em 5 conjuntos, o seu $Z$ fique acima da média das execuções. As curvas de convergência (\autoref{fig:res-conjunto1-muito-grande} a \autoref{fig:res-conjunto10-muito-grande}) ainda melhoram na última geração em todas as execuções. Com apenas 8 gerações por execução, contra cerca de 15 nas instâncias média e grande, este é o porte em que a heurística mais sente o teto de tempo reduzido, conforme antecipado no protocolo experimental. Ainda assim, o desvio padrão entre execuções foi o menor entre os portes, de 1,6\% da média.

\begin{table}[htbp]
    \caption{Função objetivo $Z$ por conjunto --- instância muito grande. Na heurística, $Z$ sem a penalidade de capacidade; $\sigma$: desvio padrão de $Z$ entre as 5 execuções do conjunto}
    \label{tab:res-z-muito-grande}
    \centering
    \begin{adjustbox}{max width=\textwidth}
    \begin{tabular}{crrrrr}
        \hline
        & \multicolumn{3}{c}{\textbf{Gurobi}} & \multicolumn{2}{c}{\textbf{Heurística (5 execuções)}} \\
        \cline{2-4} \cline{5-6}
        \textbf{Conjunto} & \textbf{Limite superior} & \textbf{Limite inferior} & \textbf{Gap} & \textbf{Melhor execução} & \textbf{Média $\pm$ $\sigma$} \\
        \hline
        01 & 3\,273\,725,0 & 0,0 & 100,0\% & 189\,250,0 & 191\,180,0 $\pm$ 2\,166,5 \\
        02 & 3\,276\,950,0 & 0,0 & 100,0\% & 217\,650,0 & 217\,570,0 $\pm$ 5\,163,9 \\
        03 & 3\,295\,100,0 & 0,0 & 100,0\% & 216\,175,0 & 212\,235,0 $\pm$ 2\,832,7 \\
        04 & 3\,289\,725,0 & 0,0 & 100,0\% & 231\,450,0 & 232\,970,0 $\pm$ 4\,856,4 \\
        05 & 3\,272\,475,0 & 0,0 & 100,0\% & 215\,075,0 & 214\,825,0 $\pm$ 2\,603,6 \\
        06 & 3\,257\,200,0 & 0,0 & 100,0\% & 181\,950,0 & 183\,880,0 $\pm$ 1\,333,2 \\
        07 & 3\,284\,350,0 & 0,0 & 100,0\% & 203\,450,0 & 202\,710,0 $\pm$ 1\,986,6 \\
        08 & 3\,261\,000,0 & 0,0 & 100,0\% & 186\,350,0 & 185\,585,0 $\pm$ 2\,232,3 \\
        09 & 3\,297\,800,0 & 0,0 & 100,0\% & 257\,875,0 & 264\,740,0 $\pm$ 7\,164,7 \\
        10 & 3\,263\,375,0 & 0,0 & 100,0\% & 207\,525,0 & 208\,250,0 $\pm$ 4\,657,9 \\
        \hline
    \end{tabular}
    \end{adjustbox}
\end{table}

\begin{figure}[htbp]
    \centering
    \begin{minipage}[b]{0.48\textwidth}
        \centering
        \includegraphics[width=\textwidth]{Figuras/resultados/muito_grande/set_01/convergence_broken_axis.png}\\
        \small (a) Convergência da heurística
    \end{minipage}
    \hfill
    \begin{minipage}[b]{0.48\textwidth}
        \centering
        \includegraphics[width=\textwidth]{Figuras/resultados/muito_grande/set_01/gurobi_vs_iterations.png}\\
        \small (b) Gurobi vs.\ execuções da heurística
    \end{minipage}

    \vspace{0.4cm}
    \begin{minipage}[b]{0.48\textwidth}
        \centering
        \includegraphics[width=\textwidth]{Figuras/resultados/muito_grande/set_01/boxplot.png}\\
        \small (c) Dispersão de $Z$ entre execuções
    \end{minipage}
    \caption{Conjunto 1 --- instância muito grande}
    \label{fig:res-conjunto1-muito-grande}
\end{figure}
\clearpage

% Conjuntos 2 a 10: desativados até as imagens estarem no Overleaf (remover \iffalse/\fi para reativar)
\iffalse
\begin{figure}[htbp]
    \centering
    \begin{minipage}[b]{0.48\textwidth}
        \centering
        \includegraphics[width=\textwidth]{Figuras/resultados/muito_grande/set_02/convergence_broken_axis.png}\\
        \small (a) Convergência da heurística
    \end{minipage}
    \hfill
    \begin{minipage}[b]{0.48\textwidth}
        \centering
        \includegraphics[width=\textwidth]{Figuras/resultados/muito_grande/set_02/gurobi_vs_iterations.png}\\
        \small (b) Gurobi vs.\ execuções da heurística
    \end{minipage}

    \vspace{0.4cm}
    \begin{minipage}[b]{0.48\textwidth}
        \centering
        \includegraphics[width=\textwidth]{Figuras/resultados/muito_grande/set_02/boxplot.png}\\
        \small (c) Dispersão de $Z$ entre execuções
    \end{minipage}
    \caption{Conjunto 2 --- instância muito grande}
    \label{fig:res-conjunto2-muito-grande}
\end{figure}
\clearpage

\begin{figure}[htbp]
    \centering
    \begin{minipage}[b]{0.48\textwidth}
        \centering
        \includegraphics[width=\textwidth]{Figuras/resultados/muito_grande/set_03/convergence_broken_axis.png}\\
        \small (a) Convergência da heurística
    \end{minipage}
    \hfill
    \begin{minipage}[b]{0.48\textwidth}
        \centering
        \includegraphics[width=\textwidth]{Figuras/resultados/muito_grande/set_03/gurobi_vs_iterations.png}\\
        \small (b) Gurobi vs.\ execuções da heurística
    \end{minipage}

    \vspace{0.4cm}
    \begin{minipage}[b]{0.48\textwidth}
        \centering
        \includegraphics[width=\textwidth]{Figuras/resultados/muito_grande/set_03/boxplot.png}\\
        \small (c) Dispersão de $Z$ entre execuções
    \end{minipage}
    \caption{Conjunto 3 --- instância muito grande}
    \label{fig:res-conjunto3-muito-grande}
\end{figure}
\clearpage

\begin{figure}[htbp]
    \centering
    \begin{minipage}[b]{0.48\textwidth}
        \centering
        \includegraphics[width=\textwidth]{Figuras/resultados/muito_grande/set_04/convergence_broken_axis.png}\\
        \small (a) Convergência da heurística
    \end{minipage}
    \hfill
    \begin{minipage}[b]{0.48\textwidth}
        \centering
        \includegraphics[width=\textwidth]{Figuras/resultados/muito_grande/set_04/gurobi_vs_iterations.png}\\
        \small (b) Gurobi vs.\ execuções da heurística
    \end{minipage}

    \vspace{0.4cm}
    \begin{minipage}[b]{0.48\textwidth}
        \centering
        \includegraphics[width=\textwidth]{Figuras/resultados/muito_grande/set_04/boxplot.png}\\
        \small (c) Dispersão de $Z$ entre execuções
    \end{minipage}
    \caption{Conjunto 4 --- instância muito grande}
    \label{fig:res-conjunto4-muito-grande}
\end{figure}
\clearpage

\begin{figure}[htbp]
    \centering
    \begin{minipage}[b]{0.48\textwidth}
        \centering
        \includegraphics[width=\textwidth]{Figuras/resultados/muito_grande/set_05/convergence_broken_axis.png}\\
        \small (a) Convergência da heurística
    \end{minipage}
    \hfill
    \begin{minipage}[b]{0.48\textwidth}
        \centering
        \includegraphics[width=\textwidth]{Figuras/resultados/muito_grande/set_05/gurobi_vs_iterations.png}\\
        \small (b) Gurobi vs.\ execuções da heurística
    \end{minipage}

    \vspace{0.4cm}
    \begin{minipage}[b]{0.48\textwidth}
        \centering
        \includegraphics[width=\textwidth]{Figuras/resultados/muito_grande/set_05/boxplot.png}\\
        \small (c) Dispersão de $Z$ entre execuções
    \end{minipage}
    \caption{Conjunto 5 --- instância muito grande}
    \label{fig:res-conjunto5-muito-grande}
\end{figure}
\clearpage

\begin{figure}[htbp]
    \centering
    \begin{minipage}[b]{0.48\textwidth}
        \centering
        \includegraphics[width=\textwidth]{Figuras/resultados/muito_grande/set_06/convergence_broken_axis.png}\\
        \small (a) Convergência da heurística
    \end{minipage}
    \hfill
    \begin{minipage}[b]{0.48\textwidth}
        \centering
        \includegraphics[width=\textwidth]{Figuras/resultados/muito_grande/set_06/gurobi_vs_iterations.png}\\
        \small (b) Gurobi vs.\ execuções da heurística
    \end{minipage}

    \vspace{0.4cm}
    \begin{minipage}[b]{0.48\textwidth}
        \centering
        \includegraphics[width=\textwidth]{Figuras/resultados/muito_grande/set_06/boxplot.png}\\
        \small (c) Dispersão de $Z$ entre execuções
    \end{minipage}
    \caption{Conjunto 6 --- instância muito grande}
    \label{fig:res-conjunto6-muito-grande}
\end{figure}
\clearpage

\begin{figure}[htbp]
    \centering
    \begin{minipage}[b]{0.48\textwidth}
        \centering
        \includegraphics[width=\textwidth]{Figuras/resultados/muito_grande/set_07/convergence_broken_axis.png}\\
        \small (a) Convergência da heurística
    \end{minipage}
    \hfill
    \begin{minipage}[b]{0.48\textwidth}
        \centering
        \includegraphics[width=\textwidth]{Figuras/resultados/muito_grande/set_07/gurobi_vs_iterations.png}\\
        \small (b) Gurobi vs.\ execuções da heurística
    \end{minipage}

    \vspace{0.4cm}
    \begin{minipage}[b]{0.48\textwidth}
        \centering
        \includegraphics[width=\textwidth]{Figuras/resultados/muito_grande/set_07/boxplot.png}\\
        \small (c) Dispersão de $Z$ entre execuções
    \end{minipage}
    \caption{Conjunto 7 --- instância muito grande}
    \label{fig:res-conjunto7-muito-grande}
\end{figure}
\clearpage

\begin{figure}[htbp]
    \centering
    \begin{minipage}[b]{0.48\textwidth}
        \centering
        \includegraphics[width=\textwidth]{Figuras/resultados/muito_grande/set_08/convergence_broken_axis.png}\\
        \small (a) Convergência da heurística
    \end{minipage}
    \hfill
    \begin{minipage}[b]{0.48\textwidth}
        \centering
        \includegraphics[width=\textwidth]{Figuras/resultados/muito_grande/set_08/gurobi_vs_iterations.png}\\
        \small (b) Gurobi vs.\ execuções da heurística
    \end{minipage}

    \vspace{0.4cm}
    \begin{minipage}[b]{0.48\textwidth}
        \centering
        \includegraphics[width=\textwidth]{Figuras/resultados/muito_grande/set_08/boxplot.png}\\
        \small (c) Dispersão de $Z$ entre execuções
    \end{minipage}
    \caption{Conjunto 8 --- instância muito grande}
    \label{fig:res-conjunto8-muito-grande}
\end{figure}
\clearpage

\begin{figure}[htbp]
    \centering
    \begin{minipage}[b]{0.48\textwidth}
        \centering
        \includegraphics[width=\textwidth]{Figuras/resultados/muito_grande/set_09/convergence_broken_axis.png}\\
        \small (a) Convergência da heurística
    \end{minipage}
    \hfill
    \begin{minipage}[b]{0.48\textwidth}
        \centering
        \includegraphics[width=\textwidth]{Figuras/resultados/muito_grande/set_09/gurobi_vs_iterations.png}\\
        \small (b) Gurobi vs.\ execuções da heurística
    \end{minipage}

    \vspace{0.4cm}
    \begin{minipage}[b]{0.48\textwidth}
        \centering
        \includegraphics[width=\textwidth]{Figuras/resultados/muito_grande/set_09/boxplot.png}\\
        \small (c) Dispersão de $Z$ entre execuções
    \end{minipage}
    \caption{Conjunto 9 --- instância muito grande}
    \label{fig:res-conjunto9-muito-grande}
\end{figure}
\clearpage

\begin{figure}[htbp]
    \centering
    \begin{minipage}[b]{0.48\textwidth}
        \centering
        \includegraphics[width=\textwidth]{Figuras/resultados/muito_grande/set_10/convergence_broken_axis.png}\\
        \small (a) Convergência da heurística
    \end{minipage}
    \hfill
    \begin{minipage}[b]{0.48\textwidth}
        \centering
        \includegraphics[width=\textwidth]{Figuras/resultados/muito_grande/set_10/gurobi_vs_iterations.png}\\
        \small (b) Gurobi vs.\ execuções da heurística
    \end{minipage}

    \vspace{0.4cm}
    \begin{minipage}[b]{0.48\textwidth}
        \centering
        \includegraphics[width=\textwidth]{Figuras/resultados/muito_grande/set_10/boxplot.png}\\
        \small (c) Dispersão de $Z$ entre execuções
    \end{minipage}
    \caption{Conjunto 10 --- instância muito grande}
    \label{fig:res-conjunto10-muito-grande}
\end{figure}
\clearpage
\fi

% ---------------------------------------------------------------------
\subsection{Viabilidade das soluções heurísticas}
\label{sec:res-viabilidade}

Como a heurística trata a capacidade como restrição flexível, a viabilidade de suas soluções precisa ser analisada à parte. A \autoref{tab:res-viabilidade} resume, por porte, quantas execuções encontraram ao menos uma solução viável e a violação de capacidade da melhor solução de cada conjunto, medida em pacientes-dia ($V_{cap}$, \autoref{eq:violacao-capacidade}) e em relação à carga de internação ($\rho_{cap}$, \autoref{eq:taxa-violacao}).

Nas instâncias pequena e média, todas as execuções encontraram soluções viáveis, e a melhor solução de todos os conjuntos respeita a capacidade. A partir da instância grande, a viabilidade torna-se exceção: apenas 5 das 100 execuções encontraram uma solução viável, e a melhor solução violou a capacidade em 7 dos 10 conjuntos. Na instância muito grande, nenhuma execução foi viável, e a melhor solução violou a capacidade nos 10 conjuntos.

A violação, contudo, é residual. Na instância grande, a melhor solução excede a capacidade em menos de um paciente-dia, em média (0,9), o que corresponde a 0,014\% da carga de internação. Na muito grande, são 24,4 pacientes-dia, ou 0,095\% da carga. Em termos práticos, trata-se de poucos pacientes que, em algum dia, seriam alocados a um quarto já lotado, enquanto mais de 99,9\% da demanda de internação é atendida dentro da capacidade. Essas soluções não podem ser aplicadas sem ajustes, mas o excesso é pequeno diante do porte das instâncias. As soluções do Gurobi, por outro lado, são sempre viáveis, pois a capacidade é uma restrição rígida do modelo exato.

\begin{table}[htbp]
    \caption{Viabilidade das soluções heurísticas por porte (melhor solução de cada conjunto, média $\pm$ desvio padrão entre os 10 conjuntos)}
    \label{tab:res-viabilidade}
    \centering
    \begin{adjustbox}{max width=\textwidth}
    \begin{tabular}{l*{4}{>{\centering\arraybackslash}p{3.3cm}}}
        \hline
        \textbf{Porte} & \textbf{Execuções que encontraram solução viável} & \textbf{Nº Violações de capacidade ($V_{cap}$)} & \textbf{Carga de internação ($N_{pd}$)} & \textbf{Taxa de violação de capacidade ($\rho_{cap} = V_{cap} / N_{pd}$)} \\
        \hline
        Pequena & 100/100 & 0,0 $\pm$ 0,0 & 397,8 $\pm$ 11,0 & 0,000\% $\pm$ 0,000\% \\
        Média & 100/100 & 0,0 $\pm$ 0,0 & 1\,611,6 $\pm$ 21,1 & 0,000\% $\pm$ 0,000\% \\
        Grande & 5/100 & 0,9 $\pm$ 0,7 & 6\,420,6 $\pm$ 41,7 & 0,014\% $\pm$ 0,011\% \\
        Muito grande & 0/50 & 24,4 $\pm$ 2,5 & 25\,601,0 $\pm$ 107,7 & 0,095\% $\pm$ 0,010\% \\
        \hline
    \end{tabular}
    \end{adjustbox}
\end{table}

% ---------------------------------------------------------------------
\subsection{Esforço computacional}
\label{sec:res-esforco}

A \autoref{tab:res-esforco} compara o esforço de busca das duas abordagens dentro do mesmo teto de tempo, medido pelo número de iterações do simplex no Gurobi e pelo número de avaliações da função objetivo na heurística. As duas medidas não têm a mesma unidade, mas a forma como variam com o porte é comparável.

O esforço do Gurobi não acompanha o crescimento das instâncias. Com o modelo cada vez maior, cada iteração se torna mais custosa e, na instância muito grande, o solver realizou apenas 1\,334 iterações em 900 segundos, sem concluir a relaxação inicial. A heurística, por sua vez, manteve uma taxa de 52 a 73 mil avaliações por segundo, chegando a 46,6 milhões de avaliações por execução na instância muito grande. Essa diferença decorre da avaliação procedural e incremental descrita na \autoref{sec:implementacao}, em que o custo de cada movimento depende apenas da internação do paciente movido, e não do tamanho do modelo.

O número de gerações, por outro lado, cai com o porte, de 25,5 na instância pequena para 8 na muito grande, pois cada geração exige refinar indivíduos com cada vez mais pacientes. Combinado com a convergência discutida nas subseções anteriores, isso explica os dois regimes observados: na instância pequena, sobram gerações e a heurística estagna; nas demais, as gerações se esgotam enquanto a busca ainda melhora.

\begin{table}[htbp]
    \caption{Esforço computacional por porte (média por execução)}
    \label{tab:res-esforco}
    \centering
    \begin{adjustbox}{max width=\textwidth}
    \begin{tabular}{lrrrr}
        \hline
        \textbf{Porte} & \textbf{Teto} & \textbf{Nº Iterações do Gurobi} & \textbf{Nº Avaliações da heurística} & \textbf{Nº Gerações} \\
        \hline
        Pequena      & 30 s  & 1\,112\,503 & 2\,202\,174  & 25,5 \\
        Média        & 90 s  & 50\,741     & 5\,841\,757  & 15,5 \\
        Grande       & 360 s & 220\,275    & 22\,175\,677 & 14,5 \\
        Muito grande & 900 s & 1\,334      & 46\,617\,002 & 8,0  \\
        \hline
    \end{tabular}
    \end{adjustbox}
\end{table}

\clearpage

% =====================================================================
\section{Discussão das métricas operacionais}
\label{sec:res-operacionais}

Enquanto a \autoref{sec:res-comparacao} compara as abordagens pelo valor agregado de $Z$, esta seção analisa o que cada solução representa para o hospital. A \autoref{tab:res-operacionais} reúne as características das instâncias e as métricas operacionais da solução do Gurobi e da melhor solução heurística de cada conjunto.

\begin{table}[htbp]
    \caption{Características das instâncias e métricas operacionais das soluções por porte (média $\pm$ desvio padrão entre os 10 conjuntos)}
    \label{tab:res-operacionais}
    \centering
    \begin{adjustbox}{max width=\textwidth}
    \begin{tabular}{lrrrr}
        \hline
        \textbf{Métrica} & \textbf{Pequena} & \textbf{Média} & \textbf{Grande} & \textbf{Muito grande} \\
        \hline
        \multicolumn{5}{l}{\textbf{Instâncias}} \\
        Horizonte (dias) & 15,1 $\pm$ 0,6 & 15,1 $\pm$ 0,3 & 15,0 $\pm$ 0,0 & 15,0 $\pm$ 0,0 \\
        Ocupação Média (todo o horizonte) & 87,9\% $\pm$ 1,9\% & 89,0\% $\pm$ 1,4\% & 89,2\% $\pm$ 0,6\% & 88,9\% $\pm$ 0,4\% \\
        Pico de Ocupação Diária & 100,0\% $\pm$ 0,0\% & 100,0\% $\pm$ 0,0\% & 100,0\% $\pm$ 0,0\% & 100,0\% $\pm$ 0,0\% \\
        Mínimo de Ocupação Diária & 11,3\% $\pm$ 6,9\% & 14,3\% $\pm$ 6,4\% & 12,3\% $\pm$ 2,6\% & 11,0\% $\pm$ 1,2\% \\
        \hline
        \multicolumn{5}{l}{\textbf{Métricas Operacionais}} \\
        \multicolumn{5}{l}{Especialidade Incorreta (contagem)} \\
        \quad Gurobi & 93,0 $\pm$ 35,7 & 369,9 $\pm$ 70,0 & 4\,618,8 $\pm$ 56,6 & 18\,470,5 $\pm$ 99,3 \\
        \quad Heurística & 93,3 $\pm$ 35,2 & 240,1 $\pm$ 54,0 & 346,2 $\pm$ 102,6 & 1\,082,6 $\pm$ 200,2 \\
        \multicolumn{5}{l}{Transferências (contagem)} \\
        \quad Gurobi & 3,7 $\pm$ 1,8 & 575,4 $\pm$ 150,3 & 4\,789,4 $\pm$ 39,9 & 19\,170,7 $\pm$ 108,2 \\
        \quad Heurística & 6,3 $\pm$ 1,9 & 52,0 $\pm$ 11,8 & 284,9 $\pm$ 30,4 & 1\,239,8 $\pm$ 88,2 \\
        \multicolumn{5}{l}{Quarto-dias Mistos (contagem)} \\
        \quad Gurobi & 11,8 $\pm$ 4,1 & 243,1 $\pm$ 38,8 & 1\,585,6 $\pm$ 19,1 & 6\,287,8 $\pm$ 29,7 \\
        \quad Heurística & 12,3 $\pm$ 4,4 & 29,3 $\pm$ 5,5 & 98,6 $\pm$ 13,8 & 539,0 $\pm$ 25,3 \\
        \multicolumn{5}{l}{Violação de Capacidade (pacientes-dia)} \\
        \quad Gurobi & 0,0 $\pm$ 0,0 & 0,0 $\pm$ 0,0 & 0,0 $\pm$ 0,0 & 0,0 $\pm$ 0,0 \\
        \quad Heurística & 0,0 $\pm$ 0,0 & 0,0 $\pm$ 0,0 & 0,9 $\pm$ 0,7 & 24,4 $\pm$ 2,5 \\
        \hline
    \end{tabular}
    \end{adjustbox}
\end{table}

\subsection{Características das instâncias}

As quatro primeiras linhas da \autoref{tab:res-operacionais} descrevem as instâncias e são iguais para qualquer solução de um mesmo conjunto, pois dependem apenas dos dados gerados. O horizonte de cerca de 15 dias e a ocupação média de 88\% a 89\% refletem a ocupação-alvo de 90\% do gerador (\autoref{sec:gerador-instancias}). O pico de 100\% em todos os conjuntos é consequência da admissão por empacotamento: os leitos ficam totalmente ocupados durante cerca de 12 dias, e a ocupação só cai nos últimos dias do horizonte, quando restam apenas as internações mais longas, o que explica o mínimo de 11\% a 14\%. Com o hospital lotado na maior parte do horizonte, não há folga de capacidade a ser explorada, cenário de alta ocupação associado, na literatura, ao bloqueio de leitos e à superlotação \cite{brandeau2004operations, he2019systematic}.

\subsection{Qualidade das alocações}

A alocação em especialidade inadequada é a principal parcela de $Z$, por ter o maior peso ($W_{spec}$) e ser contabilizada a cada dia de internação. Na instância pequena, ela responde por cerca de 85\% do $Z$ da heurística. Em relação à carga de internação, a heurística deixou 23,5\% dos pacientes-dia fora da especialidade adequada na instância pequena, 14,9\% na média, 5,4\% na grande e 4,2\% na muito grande. Boa parte desses valores é inevitável: como a especialidade de cada quarto é sorteada, a capacidade de uma especialidade pode ser menor que a sua demanda nos dias de lotação, e, na instância pequena, metade dos conjuntos não tem nenhum quarto de alguma especialidade. Contando, em cada dia, os pacientes que excedem os leitos de sua especialidade, obtém-se um mínimo de 23,2\%, 14,6\%, 4,7\% e 3,0\% dos pacientes-dia, respectivamente, valores próximos aos alcançados pela heurística. A queda com o porte reflete, portanto, a distribuição dos quartos, que se aproxima da distribuição da demanda à medida que o hospital cresce.

As transferências são a métrica em que a heurística mais se afasta de uma alocação ideal. Na instância pequena, a heurística realizou, em média, 6,3 transferências por 100 pacientes, contra 3,7 do Gurobi, e essa taxa cresce com o porte, para 13,0 na média, 17,8 na grande e 19,4 na muito grande. As transferências são criadas pelas vizinhanças parciais da busca local (N3 e N4), que as usam para contornar gargalos de capacidade nos dias de lotação, e tendem a se acumular ao longo das gerações. Parte delas é compensada na função objetivo, pois o peso de uma transferência ($W_{transf} = 50$) é metade do peso de um único paciente-dia em especialidade inadequada, de modo que transferir um paciente para o quarto correto pelo restante da internação tende a reduzir $Z$. Ainda assim, como as transferências causam desconforto ao paciente e trabalho adicional à equipe \cite{demeester2010hybrid, schafer2023combining}, esse é o principal ponto de atenção das soluções heurísticas.

Quanto à política de gênero, a heurística manteve entre 4,1\% e 8,2\% dos quartos-dia com pacientes de gêneros distintos, valor próximo ao do Gurobi na instância pequena (7,8\%). A proporção é maior na instância pequena, em que há menos quartos para separar os gêneros dentro de cada especialidade, e volta a subir na muito grande (5,6\%), em que a heurística dispõe de menos gerações. A violação de capacidade, discutida na \autoref{sec:res-viabilidade}, é nula nas soluções do Gurobi e residual nas soluções heurísticas das instâncias grande e muito grande.

Nas instâncias média, grande e muito grande, as métricas do Gurobi são muito piores que as da heurística. Na muito grande, por exemplo, 72\% dos pacientes-dia estão em especialidade inadequada e há cerca de três transferências por paciente. Esses números não refletem o modelo exato, mas a melhor solução inteira encontrada dentro do teto de tempo, ainda distante do ótimo (gap próximo de 100\%). Nesses portes, portanto, a comparação operacional deve ser lida como uma comparação entre a heurística e o que o solver consegue entregar no mesmo tempo, e não entre a heurística e a alocação ótima.

% =====================================================================
\section{Síntese}
\label{sec:res-sintese}

Os resultados mostram dois comportamentos distintos ao longo dos portes. Na instância pequena, o modelo exato é tratável: o Gurobi obteve a melhor solução em todos os conjuntos, com gaps de 1\% a 12\%, e a heurística ficou, em média, 2\% acima dele e a no máximo 7\% do ótimo. A partir da instância média, a situação se inverte: o Gurobi deixa de fechar o gap dentro do teto de tempo, e a heurística passa a obter o menor $Z$ em todos os 30 conjuntos, com valores de 54\% a 95\% menores. Esse resultado confirma, nas instâncias sintéticas deste trabalho, a dificuldade dos métodos exatos em instâncias de grande porte relatada na literatura \cite{ceschia2012modeling, demeester2010hybrid}. Em nenhum dos 40 conjuntos o Gurobi comprovou a otimalidade.

A heurística mostrou-se estável em todos os portes, com desvio padrão entre execuções de 1,6\% a 2,5\% da média, o que indica a robustez do método \cite{talbi2009metaheuristics}. Do ponto de vista operacional, as melhores soluções ficam próximas do mínimo possível de alocações em especialidade inadequada e mantêm menos de 9\% dos quartos-dia mistos, mas acumulam transferências nas instâncias maiores.

As limitações observadas são de dois tipos. Nas instâncias grande e muito grande, a melhor solução heurística pode violar a capacidade, ainda que em menos de 0,1\% da carga de internação. Além disso, o desempenho da heurística é limitado por fatores diferentes conforme o porte: na instância pequena, a busca estagna antes do fim do teto, enquanto, nas demais, o tempo termina com a busca ainda em progresso. Esses pontos indicam as direções de melhoria discutidas no capítulo seguinte: mecanismos que evitem a estagnação nas instâncias pequenas e estratégias que aumentem o número de gerações, reduzam as transferências ou eliminem a violação residual nas instâncias maiores.
